\documentclass[11pt]{article}
\usepackage[T1]{fontenc}
\usepackage{lmodern,amsmath,amssymb,amsthm,mathtools,longtable,array}
\usepackage[margin=1.05in]{geometry}
\usepackage[hidelinks]{hyperref}
\newtheorem{theorem}{Theorem}[section]
\newtheorem{lemma}[theorem]{Lemma}
\newtheorem{proposition}[theorem]{Proposition}
\newtheorem{corollary}[theorem]{Corollary}
\theoremstyle{remark}\newtheorem{remark}[theorem]{Remark}
\newcommand{\R}{\mathbb R}\newcommand{\Pconv}{\xrightarrow{\mathrm P}}
\newcommand{\norm}[1]{\left\lVert#1\right\rVert}
\newcommand{\code}[1]{{\small\nolinkurl{#1}}}
\title{Variation and spatial Simpson integration for reflected-clock\protect\\
iterated Brownian motion}
\author{Project Stardust}
\date{Reviewed editorial draft, October 29, 2026}
\begin{document}
\maketitle
\begin{abstract}
Let $Z_t=W(|B_t|)$, where $W$ and $B$ are independent standard
Brownian motions. We prove convergence in probability of oriented
spatial Simpson sums for every global $C^3$ one-form. The limit
is its ordinary spatial primitive evaluated along $Z$. The proof
uses a compact fourth-order Peano remainder and a direct quartic
tightness estimate. Separately, an exact scenery-overlap bound and
Brownian clock localization prove the uniform-grid quartic limit
$3T$, quadratic-variation divergence and the need for compensation
between raw Taylor channels.
We then prove a deterministic nonseparable $C^1$-time/$C^4$-space
chain formula, uniform control over prefixes of one common grid,
bounded random-horizon convergence for spatial one-forms, and an
explicit Bernstein realization converging in continuous-path
supremum norm. The time and space hypotheses, measurability, and
limiting topologies are stated separately. These prescribed
symmetric quadratures are not identified with a causal physical-time
divergence. The principal results have exact declarations in a
frozen Lean development; the exposition supplies full human proofs.
Novelty of this package is unresolved.
\end{abstract}

\section{The process and the quadrature prescription}
Work on a probability space carrying independent real standard
Brownian motions $B=(B_t)_{t\geq0}$ and $W=(W_s)_{s\geq0}$,
starting at zero, with continuous versions. Define
\begin{equation}\label{eq:ibm}
 Z_t=W(|B_t|),\qquad t\geq0.
\end{equation}
Thus the clock is reflected and the scenery is indexed by
nonnegative spatial time. This convention is fixed throughout.
The product Wiener realization is equivalent for these statements
to the constructed independent Gaussian scenery--clock realization
used in the formal sources. The process is jointly measurable and
has continuous paths on a common full-measure set. All samplewise
arguments below are on that set.

For a real function $g$, define the oriented spatial Simpson
increment
\begin{equation}\label{eq:simpson}
 S_g(x,y)=\frac{y-x}{6}
 \left(g(x)+4g\left(\frac{x+y}{2}\right)+g(y)\right).
\end{equation}
The sign of $y-x$ is retained. The midpoint is the spatial
midpoint of the two endpoints; it is not evaluation of $Z$ at
the midpoint of a physical time interval.

Fix $T>0$, use all uniform grids $t_i=iT/N$, $0\leq i\leq N$,
and abbreviate $Z_i=Z_{t_i}$. Define
\[
 P_{g,N}(k)=\sum_{i=0}^{k-1}S_g(Z_i,Z_{i+1}),\qquad0\leq k\leq N.
\]
The finite sums are defined directly from the integrand and the
process values. The limiting primitive will be proved to agree
with them, rather than being substituted as their definition.

\begin{theorem}\label{thm:spatial}
Let $g\in C^3(\R)$ and set
\[
 J_g(t)=\int_{Z_0}^{Z_t}g(x)\,dx,
\]
where the integral is oriented. Then $J_g$ has continuous paths,
is jointly measurable, and
\begin{equation}\label{eq:prefix}
 \max_{0\leq k\leq N}|P_{g,N}(k)-J_g(t_k)|\Pconv0.
\end{equation}
No global growth bound on $g$ or its derivatives is assumed.
\end{theorem}

The uniformity in \eqref{eq:prefix} is over prefixes of one
common grid. It is not an assertion that arbitrary partitions
or arbitrary adapted time-dependent integrands are covered.
The deterministic time--space formula and the continuous-path
realization will have their own definitions below.

Stochastic integration for iterated Brownian motion has classical
antecedents. Khoshnevisan and Lewis~\cite{kl} use a two-sided
scenery and a random partition induced by a spatial clock grid.
Our reflected-clock convention and deterministic physical-time
grids are different. Simpson-type stochastic sums have also been
studied for fractional Brownian motion~\cite{hn}; covariance or
scaling similarity does not identify those processes with
\eqref{eq:ibm}. The paper gives a verified, domain-explicit
exposition, not a claim of first construction of symmetric
integration for iterated Brownian motion.

\section{A direct quartic tightness estimate}
For an interval $[a,b]$ let $\Delta Z=Z_b-Z_a$. Conditional on
the complete clock path, the two scenery times are fixed and
$\Delta Z$ is a centered Gaussian variable of variance
\[
 L=\big||B_b|-|B_a|\big|.
\]
Independence of $W$ and $B$ is essential here. Reversal of the
order of the scenery times only changes the sign of a Brownian
increment and not its Gaussian fourth moment. Therefore
\begin{equation}\label{eq:fourthmoment}
 \mathbb E[(\Delta Z)^4\mid B]=3L^2,
 \qquad\mathbb E[(\Delta Z)^4]\leq3(b-a),
\end{equation}
using $L\leq|B_b-B_a|$ and the Brownian increment variance.
These expectations can equally be computed by Fubini on the
independent scenery--clock product space, so no unspecified
conditional version is required.

Put
\[
 Q_{4,N}=\sum_{i=0}^{N-1}|Z_{i+1}-Z_i|^4,\qquad
 M_N=\max_{i<N}|Z_{i+1}-Z_i|.
\]
Equation~\eqref{eq:fourthmoment} yields
\begin{equation}\label{eq:tight}
 \mathbb E Q_{4,N}\leq3T,\qquad
 \sup_N\mathbb P\{Q_{4,N}>A\}\leq\frac{3T}{A}
 \quad(A>0).
\end{equation}
Thus the quartic sums are tight. No independence between different
$\Delta_iZ$ is needed when summing their expectations. Continuity
of each path on $[0,T]$ gives $M_N\to0$ almost surely.

The sharper assertion $Q_{4,N}\Pconv3T$ is proved separately in
Section~\ref{sec:variation}; the quadrature proofs need only tightness.
In particular it would be incorrect to infer that assertion from
\eqref{eq:tight} or the corresponding expectation calculation.
The present argument uses only the bound it has actually proved.

\begin{lemma}\label{lem:product}
If $U_N\to0$ in probability and $V_N\geq0$ is tight, then
$U_NV_N\to0$ in probability.
\end{lemma}
\begin{proof}
For $A>0$ and $\varepsilon>0$,
\[
 \mathbb P\{|U_NV_N|>\varepsilon\}
 \leq\mathbb P\{V_N>A\}
       +\mathbb P\{|U_N|>\varepsilon/A\}.
\]
First choose $A$ to make the first term uniformly small, then
let $N\to\infty$ in the second.
\end{proof}

This elementary transfer principle is why a local remainder
modulus is sufficient; neither an almost-sure variation bound
nor global bounds on the integrand are needed.

\section{Compact Simpson remainders}
We now establish the deterministic estimate which controls
every stochastic prefix simultaneously.

\begin{lemma}\label{lem:peano}
Let $g\in C^3(\R)$ and $R>0$. There is a nonnegative modulus
$\omega_R(\delta)\to0$ as $\delta\downarrow0$ such that, for
$|x|,|y|\leq R$,
\begin{equation}\label{eq:peano}
 \left|S_g(x,y)-\int_x^y g(z)\,dz\right|
 \leq\frac13\omega_R(|y-x|)|y-x|^4.
\end{equation}
One may take $\omega_R$ to be the modulus of continuity of
$g'''$ on $[-R,R]$.
\end{lemma}
\begin{proof}
Let $d=y-x$. The case $d=0$ is immediate. On the segment between
$x$ and $y$, Taylor's theorem in integral form gives
\[
 g(x+\theta d)=g(x)+g'(x)\theta d
 +\frac{g''(x)}2\theta^2d^2+\frac{g'''(x)}6\theta^3d^3+r(\theta)
\]
for $0\leq\theta\leq1$, with
$|r(\theta)|\leq\omega_R(|d|)|d|^3/6$.
Indeed, subtract the constant $g'''(x)$ inside the third-order
integral remainder and use its uniform continuity. This proof
works for either sign of $d$.

The Simpson functional on $[0,1]$, with weights $1/6,4/6,1/6$,
integrates polynomials of degree at most three exactly. Directly,
its values on $1,\theta,\theta^2,\theta^3$ are
$1,1/2,1/3,1/4$. Thus all displayed polynomial terms cancel
against the actual integral. The Simpson average of the remainder
and its integral each have absolute value at most
$\omega_R(|d|)|d|^3/6$. Multiplying their difference by $|d|$
proves \eqref{eq:peano}.
\end{proof}

The estimate is fourth order with a vanishing compact coefficient.
It requires $g'''$ continuous, not a fourth derivative of $g$.
Equivalently, when $g=f'$, it requires a $C^4$ primitive, not a
fifth derivative. This is the regularity boundary used in the
formal compact Peano result.

For a continuous path $z:[0,T]\to\R$ with range in $[-R,R]$,
sum \eqref{eq:peano} over any grid prefix. Since the exact
oriented integrals telescope through the ordinary primitive,
\begin{equation}\label{eq:pathbound}
 \max_{k\leq N}\left|
 \sum_{i<k}S_g(z_i,z_{i+1})-\int_{z_0}^{z_k}g\right|
 \leq\frac13\omega_R(M_N(z))Q_{4,N}(z).
\end{equation}
There is one modulus and one full nonnegative quartic sum for
every prefix. No union bound with a factor growing in $N$ is
being applied.

\section{The spatial formula and measurable random horizons}
\begin{proof}[Proof of Theorem~\ref{thm:spatial}]
The function $G(x)=\int_0^xg(z)\,dz$ is $C^4$ by the fundamental
theorem of calculus, so $J_g(t)=G(Z_t)-G(Z_0)$ is continuous in
$t$ and jointly measurable. Each finite Simpson sum is measurable
because it is a finite sum of continuous functions of measurable
coordinates.

The range of a continuous sample path is bounded, so for any
$\eta>0$ choose deterministic $R$ with
$\mathbb P\{\sup_{t\leq T}|Z_t|>R\}<\eta$.
On its complementary event \eqref{eq:pathbound} applies.
There $\omega_R(M_N)\to0$ almost surely. Quartic tightness
\eqref{eq:tight} and Lemma~\ref{lem:product} show that the
right side tends to zero in probability. The probability of the
excluded event is less than $\eta$, which is arbitrary. This
proves the uniform-prefix assertion.
\end{proof}

\begin{corollary}\label{cor:random}
Let $\tau$ be any measurable random time with $0\leq\tau\leq T$.
Define
\[
 k_N=\left\lfloor\frac{N\tau}{T}\right\rfloor,\qquad
 \tau_N=\frac{T}{N}k_N.
\]
Then $P_{g,N}(k_N)\Pconv J_g(\tau)$ for $g\in C^3(\R)$.
\end{corollary}
\begin{proof}
The integer index is measurable, takes values in $\{0,\ldots,N\}$,
and obeys $0\leq\tau-\tau_N<T/N$. Evaluating the finite collection
of measurable prefix sums at that index is measurable. The error
to $J_g(\tau_N)$ is bounded by the prefix maximum in
\eqref{eq:prefix}. Pathwise continuity gives
$J_g(\tau_N)\to J_g(\tau)$ almost surely, proving the result.
\end{proof}

No stopping-time hypothesis is used. For example $\tau$ can equal
$T$ or $T/2$ according to the sign of $Z_T$. This conclusion
uses uniform approximation and continuity, not optional sampling
of a martingale.

\section{Deterministic time--space integrands}
Let $F:\R^2\to\R$ satisfy the following hypotheses: for every
fixed $t$, $F(t,\cdot)\in C^4(\R)$; for every fixed $x$,
$F(\cdot,x)\in C^1(\R)$; and the functions $\partial_tF$ and
$\partial_x^4F$ are jointly continuous. No mixed derivative or
global derivative bound is assumed. For the continuous-path
realization in Section~\ref{sec:bernstein}, we additionally state
joint continuity of $F$ explicitly, as in its formal declaration.

Define the time-frozen spatial sums
\begin{equation}\label{eq:timesums}
 P_{F,N}(k)=\sum_{i<k}
 S_{\partial_xF(t_i,\cdot)}(Z_i,Z_{i+1}),
\end{equation}
and the continuous target
\begin{equation}\label{eq:target}
 J_F(t)=F(t,Z_t)-F(0,Z_0)
       -\int_0^t\partial_tF(s,Z_s)\,ds.
\end{equation}
The target is continuous. In fact the stated hypotheses already
give joint continuity of $F$: near a fixed time $t_0$, write
$F(t,x)=F(t_0,x)+\int_{t_0}^t\partial_tF(s,x)\,ds$ and use
continuity and local boundedness of the joint time partial.
We keep the explicit hypothesis for the later source correspondence.

\begin{theorem}\label{thm:time}
Under these time--space hypotheses,
\begin{equation}\label{eq:timeprefix}
 \max_{k\leq N}|P_{F,N}(k)-J_F(t_k)|\Pconv0.
\end{equation}
In particular the full sums converge in probability to
\eqref{eq:target} at $t=T$.
\end{theorem}
\begin{proof}
On a compact rectangle $[0,T]\times[-R,R]$, joint continuity
of $\partial_x^4F$ gives one horizontal modulus uniform in $t$.
Apply Lemma~\ref{lem:peano} with $g=\partial_xF(t,\cdot)$.
The proof uses only differences of $\partial_x^4F$ at nearby
spatial points, so its modulus works for every frozen time.
Thus the maximum error in replacing \eqref{eq:timesums} by
\[
 V_{F,N}(k)=\sum_{i<k}
     \big(F(t_i,Z_{i+1})-F(t_i,Z_i)\big)
\]
tends to zero in probability, by exactly the same localization
and quartic-tightness argument as before.

It remains to identify $V_{F,N}$. The exact finite chain identity is
\begin{align}\label{eq:finitechain}
 F(t_k,Z_k)-F(0,Z_0)
 &=V_{F,N}(k)\notag\\
 &\quad+\sum_{i<k}
       \big(F(t_{i+1},Z_{i+1})-F(t_i,Z_{i+1})\big).
\end{align}
The time fundamental theorem of calculus identifies the second
sum as
\[
 \sum_{i<k}\int_{t_i}^{t_{i+1}}\partial_tF(s,Z_{i+1})\,ds.
\]
For each fixed continuous sample path, uniform continuity of
$\partial_tF$ on its compact time--space range shows that its
difference from $\int_0^{t_k}\partial_tF(s,Z_s)\,ds$ is bounded,
uniformly over $k$, by
\[
 T\,\nu_R\!\left(\sup_{|t-s|\leq T/N}|Z_t-Z_s|\right)
 \longrightarrow0.
\]
Here $\nu_R$ is a horizontal modulus of the time partial on
the compact rectangle. This is a pathwise limit. Combining it
with the spatial error and \eqref{eq:finitechain} proves
\eqref{eq:timeprefix}.
\end{proof}

The time correction is an ordinary finite Lebesgue integral on
each sample path. It is not an assumed stochastic correction
term. Joint measurability of $\partial_tF(s,Z_s)$ and integration
over a bounded interval make it a measurable random variable.
The same argument gives measurability of the finite sums, since
each frozen spatial slice has a continuous derivative.

\section{Weights and elementary checks on the prescription}
Taking $F(t,x)=h(t)G(x)$, where $h\in C^1(\R)$ and
$G'=g\in C^3(\R)$, gives
\begin{equation}\label{eq:weighted}
 \sum_{i<N}h(t_i)S_g(Z_i,Z_{i+1})\Pconv
 h(T)J_g(T)-\int_0^T h'(s)J_g(s)\,ds.
\end{equation}
The equality of this target with the one in
Theorem~\ref{thm:time} follows from cancellation of the constant
$G(Z_0)$ and $\int_0^T h'=h(T)-h(0)$.
Thus the separable time-weighted rule is contained in the full
nonseparable theorem. Random time coefficients are not included
by replacing $h$ with a random function without further work.

The spatial rule is exact for cubic one-forms. For example
\[
 S_{4z^3}(x,y)=y^4-x^4.
\]
It is not always an exact increment of a primitive: direct
polynomial expansion gives
\begin{equation}\label{eq:quintic}
 S_{5z^4}(x,y)-(y^5-x^5)=\frac{(y-x)^5}{24}.
\end{equation}
For a continuous path the sum of absolute fifth powers is at
most $M_NQ_{4,N}$, which tends to zero in probability here.
This example illustrates the finite prescription and its limiting
primitive without conflating them.

The sign in \eqref{eq:simpson} also matters. Replacing $y-x$
by $|y-x|$ would already destroy the oriented integral for a
constant one-form. Similarly, replacing the spatial midpoint
by $Z_{(t_i+t_{i+1})/2}$ changes the approximation scheme;
the deterministic Simpson polynomial cancellations no longer
apply to that expression.

\section{An explicit continuous-path realization}\label{sec:bernstein}
Under the time--space hypotheses, including explicit joint
continuity of $F$, define on $0\leq u\leq1$
\begin{equation}\label{eq:bernstein}
 A_{F,N}(u)=\sum_{k=0}^N
 \binom Nk u^k(1-u)^{N-k}P_{F,N}(k).
\end{equation}
This is an actual random polynomial assembled from the prefixes
of the one common grid. Let $j_F(u)=J_F(Tu)$.

\begin{theorem}\label{thm:path}
As random elements of $C([0,1],\R)$ with supremum norm,
\begin{equation}\label{eq:pathlimit}
 \norm{A_{F,N}-j_F}_\infty\Pconv0.
\end{equation}
The approximants and target are Borel measurable in this
continuous-path space.
\end{theorem}
\begin{proof}
The Bernstein basis is nonnegative and sums to one, so
\begin{equation}\label{eq:bernsteinbound}
 \norm{A_{F,N}-B_Nj_F}_\infty
 \leq\max_{k\leq N}|P_{F,N}(k)-j_F(k/N)|,
\end{equation}
where $B_N$ is the ordinary Bernstein approximation operator.
The right side tends to zero in probability by
Theorem~\ref{thm:time}.

For a fixed continuous $f$, Bernstein approximation gives
$\norm{B_Nf-f}_\infty\to0$. A direct
proof writes $(B_Nf)(u)=\mathbb E[f(K/N)]$ for a binomial
$K$ of parameters $N,u$. Split at $|K/N-u|\leq\delta$ and use
$\mathbb E(K/N-u)^2=u(1-u)/N\leq1/(4N)$ to obtain
\[
 \norm{B_Nf-f}_\infty
 \leq\omega_f(\delta)+\frac{2\norm f_\infty}{4N\delta^2}.
\]
First take $N\to\infty$ and then $\delta\downarrow0$.
Apply this pathwise to the continuous target $j_F$.
The resulting almost-sure Bernstein error tends to zero in
probability. Combine it with \eqref{eq:bernsteinbound}.

Each coefficient in \eqref{eq:bernstein} is measurable, hence
the finite polynomial is a measurable $C([0,1])$-valued map.
For the target, rational evaluation maps are measurable by
the measurability arguments following Theorem~\ref{thm:time}.
The Borel sigma algebra of the separable space $C([0,1])$ is
generated by these evaluations: suprema of continuous-path
distances equal suprema over a countable dense subset. Thus the
target path map is measurable as well.
\end{proof}

The approximation has $A_{F,N}(0)=0$ and
$A_{F,N}(1)=P_{F,N}(N)$, since the Bernstein basis selects the
first and last coefficients at the endpoints. At intermediate
times, all grid prefixes can enter \eqref{eq:bernstein}.
Consequently it is a polynomial smoothing, not a claimed causal
interpolation. No adaptedness or $L^p$ convergence follows merely
from \eqref{eq:pathlimit}.

\section{Exact uniform-grid quartic variation}\label{sec:variation}
The tightness estimate used above is deliberately weaker than the
library's full variation result. We now reconstruct that result
without changing the proof or domain of the Simpson formulas.
For $t_i=iT/N$, let
\[
 Q_N=\sum_{i=0}^{N-1}(Z_{i+1}-Z_i)^4,
 \qquad C_N=3\sum_{i=0}^{N-1}(|B_{t_{i+1}}|-|B_{t_i}|)^2.
\]
The second quantity is the exact scenery-conditional mean of the
first, since each scenery increment has conditional variance
$\ell_i=\big||B_{t_{i+1}}|-|B_{t_i}|\big|$.

\begin{theorem}\label{thm:quartic}
For every $T>0$, along all positive integer uniform-grid sizes,
\begin{equation}\label{eq:quarticlimit}
 Q_N\Pconv3T.
\end{equation}
\end{theorem}

The scenery increments need not be independent after conditioning
on the clock. A later traversal may revisit a previous scenery
interval. We keep all such overlaps in the variance estimate.

\subsection*{The exact Gaussian overlap calculation}
Write $a_i=|B_{t_i}|$, $b_i=|B_{t_{i+1}}|$ and
$I_i=(\min(a_i,b_i),\max(a_i,b_i)]$. Let $o_{ij}$ be the signed
overlap of these oriented intervals; its absolute value is
$|I_i\cap I_j|$. Brownian scenery increments $G_i=W_{b_i}-W_{a_i}$
satisfy, for this fixed clock sample,
\[
 \mathbb E_WG_i^2=\ell_i,\qquad
 \mathbb E_W(G_iG_j)=o_{ij}.
\]
The second equality follows by writing each increment as the
oriented integral of an interval indicator against Brownian motion.
It also follows directly by expanding the covariance
$\mathbb E(W_uW_v)=\min(u,v)$.
For centered jointly Gaussian variables with variances $l,m$ and
covariance $o$, the pair-moment formula gives
\begin{equation}\label{eq:wickquartic}
 \operatorname{Cov}(G^4,H^4)=72lm\,o^2+24o^4.
\end{equation}
To verify the constants, among pairings of four $G$'s and four
$H$'s, zero cross-pairs contribute $9l^2m^2$, two cross-pairs
contribute $\binom42^2 2!=72$ copies of $lm\,o^2$, and four
cross-pairs contribute $4!=24$ copies of $o^4$.
Subtract $\mathbb EG^4\mathbb EH^4=9l^2m^2$.
The pairing rule itself follows by expanding the Gaussian generating
function $\exp((lx^2+2oxy+my^2)/2)$ at $x^4y^4$.

If all $\ell_i\leq\delta$, then $|o_{ij}|\leq\delta$ and
\begin{equation}\label{eq:conditionalvariance}
 \mathbb E_W(Q_N-C_N)^2
 \leq96\delta^3\sum_{i,j}|o_{ij}|.
\end{equation}
Indeed, each summand in \eqref{eq:wickquartic} is bounded by its
coefficient times $\delta^3|o_{ij}|$. Every finite eighth moment
exists, so the double variance expansion is legitimate. Equivalently,
if $n(x)=\sum_i\mathbf1_{I_i}(x)$ is the return count, then
$\sum_{i,j}|o_{ij}|=\int n(x)^2dx$. Repeated visits are counted
quadratically, as required.

\subsection*{A Brownian endpoint envelope}
Define $\tau_\delta(x)=(2\delta-|x|)_+$ and
\[
 E_{N,\delta}=\sum_{i,j=0}^{N-1}
       \tau_\delta(|B_{t_i}|-|B_{t_j}|),\qquad
 p(s)=(2\pi s)^{-1/2}\quad(s>0).
\]
For two intervals starting at $a,c$ and of lengths at most
$\delta$, their overlap length is at most
$(2\delta-|a-c|)_+$. This follows by enclosing them in
$[a-\delta,a+\delta]$ and $[c-\delta,c+\delta]$ and intersecting
these two intervals. Consequently, on the event
\[
 \mathcal A_{N,\delta}
 =\{\max_i|B_{t_{i+1}}-B_{t_i}|<\delta\},
\]
the reverse triangle inequality gives $\ell_i\leq\delta$ and
$\sum_{i,j}|o_{ij}|\leq E_{N,\delta}$.

For $s<t$, condition on $B_s=x$. The later value is $x+G$, where
$G$ is independent normal with variance $t-s$. The folded variable
$|x+G|$ has density on $[0,\infty)$ bounded by $2p(t-s)$,
uniformly in $x$. The tent has height at most $2\delta$, support
length $4\delta$, and hence integral at most $8\delta^2$.
Thus
\begin{equation}\label{eq:tentmean}
 \mathbb E\tau_\delta(|B_s|-|B_t|)
 \leq16\delta^2p(t-s).
\end{equation}
The diagonal is instead exactly $2\delta$; it is not evaluated
through a singular zero-variance density formula. Set $h=T/N$.
Since $p(kh)=p(h)/\sqrt k$ and
\[
 \sum_{k=1}^{N}\frac1{\sqrt k}\leq2\sqrt N,
\]
summing the ordered off-diagonal pairs gives
\begin{equation}\label{eq:endpointenergy}
 \mathbb EE_{N,\delta}
 \leq2N\delta+64\delta^2p(h)N\sqrt N
 =2N\delta+64p(T)N^2\delta^2.
\end{equation}
The inverse-square-root sum follows by telescoping
$1/\sqrt k\leq2(\sqrt k-\sqrt{k-1})$.

Let $m_{20}=\mathbb E|G_0|^{20}$ for a standard normal $G_0$;
this is a finite fixed constant. The union bound and Markov's
inequality give
\begin{equation}\label{eq:badclock}
 \mathbb P(\mathcal A_{N,\delta}^c)
 \leq N(T/N)^{10}m_{20}\delta^{-20}.
\end{equation}
Choose $\delta=\delta_N=N^{-4/9}$. The right side is
$T^{10}m_{20}N^{-1/9}\to0$, whereas
\begin{equation}\label{eq:weightedenergy}
 \delta_N^3\mathbb EE_{N,\delta_N}
 \leq2N^{-7/9}+64p(T)N^{-2/9}\longrightarrow0.
\end{equation}
For any $\varepsilon>0$, conditional Chebyshev on the good-clock
event, followed by Fubini, now yields
\begin{equation}\label{eq:quarticfluctuation}
 \mathbb P(|Q_N-C_N|\geq\varepsilon)
 \leq\mathbb P(\mathcal A_{N,\delta_N}^c)
     +\frac{96}{\varepsilon^2}\delta_N^3\mathbb EE_{N,\delta_N}
 \longrightarrow0.
\end{equation}
On the bad event the conditional probability is bounded by one.
No global joint $L^2$ error bound is inferred by removing that event.

\subsection*{The reflected-clock correction}
For real $a,b$ define $r(a,b)=2(|ab|-ab)\geq0$. Direct expansion gives
\begin{equation}\label{eq:reflectionidentity}
 (|b|-|a|)^2=(b-a)^2-r(a,b).
\end{equation}
Let $R_N=\sum_i r(B_{t_i},B_{t_{i+1}})$.
If $|b-a|\leq\delta$, then
\begin{equation}\label{eq:penaltytent}
 r(a,b)\leq\delta\tau_\delta(a).
\end{equation}
When $ab\geq0$ the left side is zero. Otherwise the endpoints
straddle zero, so $|a|\leq\delta$ and
$r(a,b)=4|ab|\leq(b-a)^2\leq\delta^2$;
also $\tau_\delta(a)\geq\delta$. This proves the bound in both cases.
Therefore, on $\mathcal A_{N,\delta}$,
\[
 0\leq R_N\leq\delta\sum_i\tau_\delta(B_{t_i}).
\]
The initial term is bounded by $2\delta^2$. At positive grid times,
the ordinary Gaussian density and the same tent-mass bound give
$\mathbb E\tau_\delta(B_{t_i})\leq16\delta^2p(t_i)$.
Summing as before, with a harmless upper bound including one extra
positive time, gives
\[
 \mathbb E\left[\delta\sum_i\tau_\delta(B_{t_i})\right]
 \leq2\delta^2+32p(T)N\delta^3.
\]
For $\delta_N=N^{-4/9}$ this tends to zero. Markov's inequality,
with the bad-clock probability kept separate, implies $R_N\Pconv0$.

Finally $U_N=\sum_i(B_{t_{i+1}}-B_{t_i})^2$ has expectation $T$
and variance $2T^2/N$, by independent Brownian increments. Thus
$U_N\to T$ in $L^2$. Summing \eqref{eq:reflectionidentity} gives
$C_N=3U_N-3R_N\Pconv3T$. Combining this with
\eqref{eq:quarticfluctuation} proves Theorem~\ref{thm:quartic}.

\section{Divergent raw channels and finite compensated sums}
Let
\[
 V_N=\sum_i(Z_{i+1}-Z_i)^2,\qquad
 m_N=\max_i|Z_{i+1}-Z_i|.
\]
Path continuity on $[0,T]$ implies $m_N\to0$ almost surely, and
the elementary inequality
\begin{equation}\label{eq:quarticquadratic}
 Q_N\leq m_N^2V_N
\end{equation}
keeps the quartic and quadratic sums distinct.

\begin{corollary}\label{cor:quadratic}
For $T>0$ and every fixed real $K$,
\[
 \mathbb P(V_N\leq K)\longrightarrow0.
\]
Neither $V_N$ nor the raw first Taylor channel
$F_N=\sum_i2Z_i(Z_{i+1}-Z_i)$ has a finite probability limit.
Nevertheless their sum is exactly $Z_T^2-Z_0^2$ for every grid.
\end{corollary}
\begin{proof}
For $K<0$ the event is empty. For $K\geq0$, on
$\{V_N\leq K\}$, \eqref{eq:quarticquadratic} gives
$Q_N\leq Km_N^2$. Since $Q_N\Pconv3T>0$ and
$Km_N^2\Pconv0$, the event has probability tending to zero.
If $V_N\Pconv Y$ for a finite random variable, choose $K$
with $\mathbb P(|Y|>K)<1/4$. For all sufficiently large $N$,
$\mathbb P(|V_N-Y|>1)<1/4$ would give
$\mathbb P(V_N\leq K+1)>1/2$, contradicting the lower-tail result.

Expanding $(Z_{i+1})^2-Z_i^2$ and telescoping gives
\[
 F_N+V_N=Z_T^2-Z_0^2.
\]
A finite probability limit for $F_N$ would therefore give one for
$V_N$ by subtraction from this fixed finite endpoint variable.
This is impossible. The compensation is essential even for the
polynomial $z^2$; it cannot be replaced by separately convergent
raw first and second channels.
\end{proof}

Burdzy's classical \emph{Variation of iterated Brownian motion}
\cite{burdzy} already establishes deterministic quartic variation
for two-sided-scenery IBM. Khoshnevisan--Lewis \cite{kl}, Theorem 4.3,
proves a stronger weighted quartic limit on its induced random
partition. These are explicit antecedents. Our source-mapped proof
retains the reflected-clock model, all uniform physical grids and
convergence in probability; it claims no priority for the quartic
phenomenon and no stronger almost-sure or weighted limit here.

\section{What this integration result means}
The targets are genuine pathwise finite spatial or time integrals,
and the approximations converge under the actual positive
scenery--clock law. Here ``oriented'' refers to the signed spatial
increment $y-x$, not to a signed probability measure for $Z$.
This distinguishes the present setting from the higher-order
signed heat pseudoprocess calculus.

The proof is local in the spatial range. It admits unbounded
global $C^3$ one-forms and $C^4$ spatial slices because every
continuous sample path has a compact range. It provides
convergence in probability, not integrability of arbitrary
unbounded targets. Neither expectation convergence nor strong
$L^2$ convergence has been asserted without moment hypotheses.

It also does not construct a general physical-time adapted
stochastic integral. Its spatial one-form depends only on the
current spatial argument, and its time--space integrand comes
from a deterministic primitive $F$. The rule has not been
identified with the original intrinsic divergence or with an
operator acting on arbitrary predictable integrands. Bounded
random horizons are handled by approximation, not a martingale
stopping theorem. Nonsmooth It\^o--Tanaka formulas and arbitrary
partitions have their own proof obligations.

The variation and raw-channel results above are independent of the
Simpson proof, which needs only quartic tightness. Intrinsic
graph-dual obstructions remain separately scoped; no stochastic-domain
assertion is used to justify a deterministic quadrature estimate.

\section{Formal correspondence and preliminary comparisons}
The frozen mathematical baseline is
\begin{center}\small\ttfamily
bab688bb6140cd6316c80680095a914ea0d346db.
\end{center}
The namespace prefix is
\begin{center}\code{OperatorDerivativeRewrite.Foundation}\end{center}
The companion \code{theorem-map.json} contains exact declaration
links, source line numbers and hashes.

\begin{longtable}{@{}p{0.25\textwidth}p{0.69\textwidth}@{}}
\textit{Paper result}&\textit{Lean namespace and declaration}\\\hline
Compact spatial estimate&\code{IBMSimpsonPeano.compact_error_bound}\\[3pt]
Spatial convergence&\code{IBMActualSimpson.sums_tendstoInMeasure}\\[3pt]
Uniform spatial prefixes&\code{IBMUniformSimpsonProcess.maximum_tendstoInMeasure}\\[3pt]
Random horizons&\code{IBMSimpsonRandomHorizon.stopped_sums_tendstoInMeasure}\\[3pt]
Time--space compact estimate&\code{IBMTimeSpacePeano.compact_error_bound}\\[3pt]
Time quadrature&\code{IBMTimeSpaceQuadrature.increments_tendsto}\\[3pt]
Time--space chain&\code{IBMTimeSpaceSimpson.sums_tendstoInMeasure}\\[3pt]
Uniform time--space prefixes&\code{IBMUniformTimeSpace.maximum_tendstoInMeasure}\\[3pt]
Time weights&\code{IBMTimeWeightedSimpson.weighted_tendstoInMeasure}\\
Quartic variation&\code{IBMUniformQuarticVariation.quartic_tendstoInMeasure}\\[3pt]
Quadratic divergence&\code{IBMActualQuadraticDivergence.quadratic_lower_tail}\\[3pt]
Raw first channel&\code{IBMActualQuadraticDivergence.first_no_finite_limit}\\
\end{longtable}

The continuous-path statements are
\code{IBMContinuousSimpsonRealization.approximation_error_bound}
and \code{approximation_tendstoInMeasure} in the same namespace.

The formal proof uses its verified quartic-variation machinery to
control compact Peano errors. The human proof above reorganizes
the same application using the weaker direct Gaussian moment
bound \eqref{eq:tight}. That bound is proved explicitly here;
no separate named Lean declaration for this shortened argument
is claimed. The endpoints, smoothness assumptions and probability
topologies of the principal convergence statements are those
in the formal declarations. The nonseparable chain theorem in
the source does not require joint continuity of $F$ as a separate
input; the continuous-path declaration retains it explicitly.

The baseline receipt records a complete successful Lean build
and 54,943 declarations audited with only the standard three
axioms. This editorial work changes no mathematical source and
does not claim a new full build. Source hashes, baseline evidence,
compiler diagnostics and visual-review receipts accompany the
paper. The earlier signed-Simpson working note remains intact;
this paper replaces its mechanism summaries with proofs for the
focused quadrature and continuous-path scope.

The primary comparison~\cite{kl} constructs symmetric integrals
for a two-sided-scenery IBM using an induced random partition.
It is not an identification with the reflected-clock convention
or deterministic grids used here. Harnett and Nualart~\cite{hn}
study the critical fractional Brownian Simpson problem and
distributional limits. These are close antecedents for numerical
prescriptions, but their process and limiting regime differ.
The compact-remainder method itself is elementary; this paper
does not turn its formal verification into a novelty claim.
Priority for the combined domain, prefix and path-space package
requires further targeted comparison.

\begin{thebibliography}{9}
\bibitem{burdzy} K. Burdzy,
Variation of iterated Brownian motion,
\emph{CRM Proceedings and Lecture Notes} \textbf{5} (1994), 35--53.
\href{https://hdl.handle.net/1773/2252}{University of Washington record}.
\bibitem{kl} D. Khoshnevisan and T. M. Lewis,
Stochastic calculus for Brownian motion on a Brownian fracture,
\emph{Ann. Appl. Probab.} \textbf{9} (1999), 629--667.
\url{https://www.math.utah.edu/~davar/ps-pdf-files/ibm.pdf}.
\bibitem{hn} D. Harnett and D. Nualart,
On Simpson's rule and fractional Brownian motion with $H=1/10$,
\emph{J. Theor. Probab.} \textbf{28} (2015), 1651--1688.
\href{https://doi.org/10.1007/s10959-014-0552-1}{doi:10.1007/s10959-014-0552-1}.
\end{thebibliography}
\end{document}
